% GENERATED FILE. Edit canonical lessons, then run npm run generate:books.
% canonical-source-hash: fnv1a64:428222c83081dc2c
% canonical-lessons: TE-C497-koppadu, TE-C497-mandipadu, TE-C497-kasuruko, TE-C497-avesam, TE-C497-ciraku

\chapter{Losing Your Temper}
\label{ch:te-losing-your-temper}
\hlchaptermodality{497}

\begin{chapteropening}
\textbf{By the end of this chapter:} \emph{I can say that someone got angry, fumed or snapped.}

Complete the last lesson of chapter 497: I can say that someone got angry, fumed or snapped.
\end{chapteropening}

\section[\texorpdfstring{\te{కోప్పడు} (kōppaḍu)}{కోప్పడు (kōppaḍu)}]{\te{కోప్పడు} (kōppaḍu) — to get angry}

\label{lesson:TE-C497-koppadu}



\begin{quote}
Before the new one: say the Telugu for to grieve inwardly, then the Telugu for despair.
\end{quote}

\subsection*{You'll want to know: \te{కోప్పడు}}
\textbf{\te{కోప్పడు}} — \emph{kōppaḍu} — \textquotedblleft{}to get angry\textquotedblright{}.

\begin{quote}
\te{ఆలస్యంగా} \te{వచ్చినందుకు} \te{నాన్న} \te{నా} \te{మీద} \te{కోప్పడ్డారు}.
\end{quote}

\begin{quote}
\emph{ālasyaṅgā vaccinanduku nānna nā mīda kōppaḍḍāru.}
\end{quote}

\begin{quote}
Dad got angry with me for coming late.
\end{quote}

Past (I …): \te{కోప్పడ్డాను}, \emph{kōppaḍḍānu}.

The person you are angry with takes \te{మీద}: \te{నా} \te{మీద} \te{కోప్పడ్డారు}.

\begin{cousinweb}
\textbf{\te{కోప్పడు}} (\emph{kōppaḍu}) is built from \te{కోపం}, \textquotedblleft{}anger\textquotedblright{}, from Sanskrit \emph{kōpa} (CDIAL 3516), and the native verb \emph{paḍu}, \textquotedblleft{}to feel, to suffer\textquotedblright{} (DEDR 3853); the two have fused into one word with a doubled \emph{p}.
\end{cousinweb}

\begin{grammarlens}[title={What you've built}]
One more word for describing what you lived and felt.
\end{grammarlens}

\subsection*{Your turn}
\begin{itemize}
\raggedright
  \item \emph{Say it:} \emph{kōppaḍu}
  \item \emph{Say it:} \emph{kōppaḍu}, once more
  \item \emph{Say it:} \emph{ālasyaṅgā vaccinanduku nānna nā mīda kōppaḍḍāru.} — Dad got angry with me for coming late.
  \item \emph{Recall:} say the Telugu for to grieve inwardly, then the Telugu for despair, then say \emph{kōppaḍu} again
\end{itemize}

\begin{tcolorbox}[breakable,colback=teal!4,colframe=teal!35!black,title={Before you move on}]
What does \te{కోప్పడు} mean? (\textquotedblleft{}To get angry\textquotedblright{}.) Say it once more.
\end{tcolorbox}
\section[\texorpdfstring{\te{మండిపడు} (maṇḍipaḍu)}{మండిపడు (maṇḍipaḍu)}]{\te{మండిపడు} (maṇḍipaḍu) — to flare up; to fume}

\label{lesson:TE-C497-mandipadu}



\begin{quote}
Before the new one: say the Telugu for despair, then the Telugu for to get angry.
\end{quote}

\subsection*{You'll want to know: \te{మండిపడు}}
\textbf{\te{మండిపడు}} — \emph{maṇḍipaḍu} — \textquotedblleft{}to flare up; to fume\textquotedblright{}.

\begin{quote}
\te{బిల్లు} \te{చూసి} \te{నాన్న} \te{మండిపడ్డారు}.
\end{quote}

\begin{quote}
\emph{billu cūsi nānna maṇḍipaḍḍāru.}
\end{quote}

\begin{quote}
Seeing the bill, dad flared up.
\end{quote}

Past (I …): \te{మండిపడ్డాను}, \emph{maṇḍipaḍḍānu}.

Literally to blaze: \te{మండు} is to burn.

\begin{cousinweb}
\textbf{\te{మండిపడు}} (\emph{maṇḍipaḍu}) is native Dravidian: DEDR 4680 lists Telugu \emph{maṇḍu}, \textquotedblleft{}to burn, to blaze, to be in a fury\textquotedblright{}, with Tamil \emph{maṇṭu}, \textquotedblleft{}to blaze up\textquotedblright{}; the native verb \emph{paḍu}, \textquotedblleft{}to feel, to suffer\textquotedblright{} (DEDR 3853) completes the verb.
\end{cousinweb}

\begin{grammarlens}[title={What you've built}]
One more word for describing what you lived and felt.
\end{grammarlens}

\subsection*{Your turn}
\begin{itemize}
\raggedright
  \item \emph{Say it:} \emph{maṇḍipaḍu}
  \item \emph{Say it:} \emph{maṇḍipaḍu}, once more
  \item \emph{Say it:} \emph{billu cūsi nānna maṇḍipaḍḍāru.} — Seeing the bill, dad flared up.
  \item \emph{Recall:} say the Telugu for despair, then the Telugu for to get angry, then say \emph{maṇḍipaḍu} again
\end{itemize}

\begin{tcolorbox}[breakable,colback=teal!4,colframe=teal!35!black,title={Before you move on}]
What does \te{మండిపడు} mean? (\textquotedblleft{}To flare up; to fume\textquotedblright{}.) Say it once more.
\end{tcolorbox}
\section[\texorpdfstring{\te{కసురుకో} (kasurukō)}{కసురుకో (kasurukō)}]{\te{కసురుకో} (kasurukō) — to snap at}

\label{lesson:TE-C497-kasuruko}



\begin{quote}
Before the new one: say the Telugu for to get angry, then the Telugu for to flare up.
\end{quote}

\subsection*{You'll want to know: \te{కసురుకో}}
\textbf{\te{కసురుకో}} — \emph{kasurukō} — \textquotedblleft{}to snap at\textquotedblright{}.

\begin{quote}
\te{ఆకలితో} \te{ఉన్న} \te{అన్నయ్య} \te{అందరి} \te{మీదా} \te{కసురుకున్నాడు}.
\end{quote}

\begin{quote}
\emph{ākalitō unna annayya andari mīdā kasurukunnāḍu.}
\end{quote}

\begin{quote}
Hungry, my brother snapped at everyone.
\end{quote}

Past (I …): \te{కసురుకున్నాను}, \emph{kasurukunnānu}.

Takes \te{మీద}, like \te{కోప్పడు}: \te{అందరి} \te{మీదా} \te{కసురుకున్నాడు}.

\begin{cousinweb}
\textbf{\te{కసురుకో}} (\emph{kasurukō}) is native Dravidian: DEDR 1089 lists Telugu \emph{kasaru}, \textquotedblleft{}to scold, to rebuke; anger\textquotedblright{}, with Malayalam \emph{kayakka}, \textquotedblleft{}to quarrel\textquotedblright{}. The ending -\te{కో} makes the verb reflexive.
\end{cousinweb}

\begin{grammarlens}[title={What you've built}]
One more word for describing what you lived and felt.
\end{grammarlens}

\subsection*{Your turn}
\begin{itemize}
\raggedright
  \item \emph{Say it:} \emph{kasurukō}
  \item \emph{Say it:} \emph{kasurukō}, once more
  \item \emph{Say it:} \emph{ākalitō unna annayya andari mīdā kasurukunnāḍu.} — Hungry, my brother snapped at everyone.
  \item \emph{Recall:} say the Telugu for to get angry, then the Telugu for to flare up, then say \emph{kasurukō} again
\end{itemize}

\begin{tcolorbox}[breakable,colback=teal!4,colframe=teal!35!black,title={Before you move on}]
What does \te{కసురుకో} mean? (\textquotedblleft{}To snap at\textquotedblright{}.) Say it once more.
\end{tcolorbox}
\section[\texorpdfstring{\te{ఆవేశం} (āvēśaṁ)}{ఆవేశం (āvēśaṁ)}]{\te{ఆవేశం} (āvēśaṁ) — fury; passion}

\label{lesson:TE-C497-avesam}



\begin{quote}
Before the new one: say the Telugu for to flare up, then the Telugu for to snap at.
\end{quote}

\subsection*{You'll want to know: \te{ఆవేశం}}
\textbf{\te{ఆవేశం}} — \emph{āvēśaṁ} — \textquotedblleft{}fury; passion\textquotedblright{}.

\begin{quote}
\te{ఆవేశంలో} \te{అతను} \te{తలుపు} \te{గట్టిగా} \te{వేశాడు}.
\end{quote}

\begin{quote}
\emph{āvēśaṁlō atanu talupu gaṭṭigā vēśāḍu.}
\end{quote}

\begin{quote}
In a fury he slammed the door.
\end{quote}

\te{ఆవేశంలో}, in the heat of the moment; \te{ఆవేశంగా}, passionately.

\begin{cousinweb}
\textbf{\te{ఆవేశం}} (\emph{āvēśaṁ}) is borrowed from Sanskrit \emph{āveśa}, \textquotedblleft{}entering, taking possession; frenzy, wrath\textquotedblright{} (Monier-Williams): a feeling that takes hold of you.
\end{cousinweb}

\begin{grammarlens}[title={What you've built}]
One more word for describing what you lived and felt.
\end{grammarlens}

\subsection*{Your turn}
\begin{itemize}
\raggedright
  \item \emph{Say it:} \emph{āvēśaṁ}
  \item \emph{Say it:} \emph{āvēśaṁ}, once more
  \item \emph{Say it:} \emph{āvēśaṁlō atanu talupu gaṭṭigā vēśāḍu.} — In a fury he slammed the door.
  \item \emph{Recall:} say the Telugu for to flare up, then the Telugu for to snap at, then say \emph{āvēśaṁ} again
\end{itemize}

\begin{tcolorbox}[breakable,colback=teal!4,colframe=teal!35!black,title={Before you move on}]
What does \te{ఆవేశం} mean? (\textquotedblleft{}Fury; passion\textquotedblright{}.) Say it once more.
\end{tcolorbox}
\section[\texorpdfstring{\te{చిరాకు} (cirāku)}{చిరాకు (cirāku)}]{\te{చిరాకు} (cirāku) — irritation; annoyance}

\label{lesson:TE-C497-ciraku}



\begin{quote}
Before the new one: say the Telugu for to snap at, then the Telugu for fury.
\end{quote}

\subsection*{You'll want to know: \te{చిరాకు}}
\textbf{\te{చిరాకు}} — \emph{cirāku} — \textquotedblleft{}irritation; annoyance\textquotedblright{}.

\begin{quote}
\te{ఎండ} \te{వల్ల} \te{పిల్లలకి} \te{చిరాకుగా} \te{ఉంది}.
\end{quote}

\begin{quote}
\emph{eṇḍa valla pillalaki cirākugā undi.}
\end{quote}

\begin{quote}
The heat is making the children cranky.
\end{quote}

With the dative: \te{నాకు} \te{చిరాకుగా} \te{ఉంది}, I feel irritable.

\begin{cousinweb}
\textbf{\te{చిరాకు}} (\emph{cirāku}) is native Dravidian: DEDR 1961 lists Telugu \emph{cirāku}, \textquotedblleft{}anger, displeasure, crossness\textquotedblright{}, beside Tamil \emph{ciraṟu}, \textquotedblleft{}to sulk, to disagree\textquotedblright{}.
\end{cousinweb}

\begin{grammarlens}[title={What you've built}]
One more word for describing what you lived and felt.
\end{grammarlens}

\subsection*{Your turn}
\begin{itemize}
\raggedright
  \item \emph{Say it:} \emph{cirāku}
  \item \emph{Say it:} \emph{cirāku}, once more
  \item \emph{Say it:} \emph{eṇḍa valla pillalaki cirākugā undi.} — The heat is making the children cranky.
  \item \emph{Recall:} say the Telugu for to snap at, then the Telugu for fury, then say \emph{cirāku} again
\end{itemize}

\begin{tcolorbox}[breakable,colback=teal!4,colframe=teal!35!black,title={Before you move on}]
What does \te{చిరాకు} mean? (\textquotedblleft{}Irritation; annoyance\textquotedblright{}.) Say it once more.
\end{tcolorbox}
